% chapter.tex (Chapter 06)
\chapter{ডেটাবেজ ম্যানেজমেন্ট সিস্টেম}

\begin{learningbox}
এই অধ্যায়ে ডেটা, ডেটাবেজ, DBMS, RDBMS, টেবিল, ফিল্ড, রেকর্ড, কী, রিলেশন, কুয়েরি, সর্টিং, ইনডেক্সিং, ফর্ম, রিপোর্ট, কর্পোরেট ও সরকারি ডেটাবেজ, ডেটা নিরাপত্তা এবং এনক্রিপশন ধারাবাহিকভাবে শেখানো হয়েছে। HSC পরীক্ষায় এই অধ্যায় থেকে সংজ্ঞা, পার্থক্য, প্রয়োগ, SQL কুয়েরি এবং ব্যবহারিক কাজ খুব বেশি গুরুত্বপূর্ণ।
\end{learningbox}

\begin{keypointbox}{অধ্যায়ের ধারণা-মানচিত্র}
\begin{center}
\begin{tabular}{|p{0.22\textwidth}|p{0.68\textwidth}|}
\hline
ভিত্তি & ডেটা, ইনফরমেশন, ফিল্ড, রেকর্ড, টেবিল, ডেটাবেজ \\
\hline
ব্যবস্থাপনা & DBMS, RDBMS, DBA, ডেটা টাইপ, ডেটাবেজ তৈরি \\
\hline
উদ্ধার ও বিশ্লেষণ & Query, SQL, Sorting, Indexing, Relation \\
\hline
উপস্থাপন & Form, Report, Corporate/Government database \\
\hline
সুরক্ষা & Access control, Backup, Encryption, Data security \\
\hline
\end{tabular}
\end{center}
\end{keypointbox}

\begin{boardbox}{বোর্ড ফোকাস}
প্রায়ই আসা বিষয়: DBMS-এর সুবিধা, DBMS ও RDBMS-এর পার্থক্য, Primary key ও Foreign key, One-to-many relation, Query-এর কাজ, Sorting ও Indexing-এর পার্থক্য, Form ও Report-এর ব্যবহার, ডেটা নিরাপত্তা এবং এনক্রিপশন।
\end{boardbox}
% Chapter 06 overview
\section{ডেটাবেজ ম্যানেজমেন্ট সিস্টেম — পরিচিতি}
\begin{itemize}
\item DBMS কি এবং মৌলিক উদ্দেশ্য
\end{itemize}

% ২. ডেটাবেজের মৌলিক উপাদান
\section{ডেটাবেজের মৌলিক উপাদান}

\subsection{Field}
\begin{itemize}
\item Field কী, উদাহরণ: Roll, Name, Address, Date of Birth
\end{itemize}

\subsection{Record, Table, File}
\begin{itemize}
\item Record কী, Table-এর row ও column, File সংজ্ঞা
\end{itemize}

\subsection{Entity, Attribute, Tuple, Domain}
\begin{itemize}
\item Entity ও attribute, Tuple (relational record), Domain বিষয়ক ব্যাখ্যা
\end{itemize}

% ৩. ডেটাবেজ ম্যানেজমেন্ট সিস্টেম (DBMS)
\section{ডেটাবেজ ম্যানেজমেন্ট সিস্টেম}

\subsection{DBMS কী}
\begin{itemize}
\item DBMS সংজ্ঞা: database ও user-এর মধ্যে মধ্যস্থ সফটওয়্যার
\end{itemize}

\subsection{DBMS-এর কাজ}
\begin{itemize}
\item Database তৈরি, Table তৈরি, Insert/Update/Delete, Query, Report
\item Sorting, Indexing, Security, Backup, Recovery, Integrity
\end{itemize}

\subsection{সুবিধা ও অসুবিধা}
\begin{itemize}
\item সুবিধা: redundancy কমা, consistency, multi-user support
\item অসুবিধা: খরচ, skilled user প্রয়োজন, system failure-এর ঝুঁকি
\end{itemize}

% ৪. RDBMS
\section{Relational Database Management System (RDBMS)}

\subsection{RDBMS কী}
\begin{itemize}
\item Table-based storage, relation between tables, primary key
\end{itemize}

\subsection{বৈশিষ্ট্য ও ব্যবহার}
\begin{itemize}
\item Characteristics: table, keys, integrity, query language, multi-user
\item Uses: banking, education, hospital, reservation systems, e-commerce
\end{itemize}

\subsection{জনপ্রিয় সফটওয়্যার}
\begin{itemize}
\item Oracle, MySQL, SQL Server, PostgreSQL, MS Access, DB2, Sybase
\end{itemize}

% ৫. ডেটাবেজ তৈরি
\section{ডেটাবেজ তৈরি}
ভালো ডেটাবেজ তৈরির আগে কী তথ্য রাখা হবে, কোন কোন টেবিল লাগবে, কোন ফিল্ডের ডেটা টাইপ কী হবে এবং কোন ফিল্ডটি রেকর্ডকে অনন্যভাবে শনাক্ত করবে তা ঠিক করতে হয়।

\subsection{তৈরির ধাপ}
\begin{keypointbox}{ডেটাবেজ তৈরির ধাপ}
\begin{enumerate}
    \item উদ্দেশ্য নির্ধারণ: ডেটাবেজটি কোন কাজ করবে।
    \item এন্টিটি নির্বাচন: Student, Teacher, Product, Sales ইত্যাদি।
    \item ফিল্ড নির্ধারণ: Roll, Name, Class, Mobile, GPA ইত্যাদি।
    \item ডেটা টাইপ নির্ধারণ: Text, Number, Date/Time ইত্যাদি।
    \item Primary key নির্বাচন: যেমন StudentID বা Roll।
    \item টেবিল তৈরি, ডেটা এন্ট্রি, সম্পর্ক তৈরি এবং পরীক্ষা।
\end{enumerate}
\end{keypointbox}

\subsection{Table ও Field ডিজাইন}
\begin{examplebox}{Student ডেটাবেজ ডিজাইন}
\begin{center}
\begin{tabular}{|l|l|l|}
\hline
Field Name & Data Type & Note \\
\hline
StudentID & Number/AutoNumber & Primary Key \\
Name & Short Text & ছাত্রের নাম \\
GroupName & Short Text & বিভাগ \\
BirthDate & Date/Time & জন্মতারিখ \\
Mobile & Short Text & শুরুতে 0 থাকায় Text ভালো \\
\hline
\end{tabular}
\end{center}
\end{examplebox}

\begin{mistakebox}{সাধারণ ভুল}
মোবাইল নম্বর বা রোল নম্বর সব সময় গণনার জন্য ব্যবহৃত হয় না। তাই অনেক ক্ষেত্রে এগুলো Number না রেখে Text রাখা নিরাপদ, বিশেষ করে শুরুতে 0 থাকতে পারে এমন ডেটার ক্ষেত্রে।
\end{mistakebox}

\begin{practicebox}
একটি লাইব্রেরি ডেটাবেজের জন্য Book, Member ও Issue টেবিলের সম্ভাব্য ফিল্ড এবং Primary key নির্ধারণ করো।
\end{practicebox}

% ৬. ডেটা টাইপ
\section{ডেটা টাইপ}
ডেটা টাইপ নির্ধারণ করে একটি ফিল্ডে কী ধরনের মান রাখা যাবে। সঠিক ডেটা টাইপ ব্যবহার করলে স্টোরেজ সাশ্রয় হয়, ভুল ডেটা এন্ট্রি কমে এবং কুয়েরি/গণনা সহজ হয়।

\subsection{প্রয়োগ}
\begin{center}
\begin{tabular}{|p{0.22\textwidth}|p{0.46\textwidth}|p{0.20\textwidth}|}
\hline
ডেটা টাইপ & ব্যবহার & উদাহরণ \\
\hline
Short Text & ছোট লেখা, কোড, মোবাইল নম্বর & Name, Mobile \\
\hline
Long Text/Memo & বড় বিবরণ বা মন্তব্য & Address, Note \\
\hline
Number & গাণিতিক হিসাবযোগ্য সংখ্যা & Marks, Quantity \\
\hline
Date/Time & তারিখ ও সময় & BirthDate \\
\hline
Currency & টাকা-পয়সা & Salary, Price \\
\hline
AutoNumber & স্বয়ংক্রিয় অনন্য নম্বর & StudentID \\
\hline
Yes/No & সত্য/মিথ্যা বা হ্যাঁ/না মান & Paid, Active \\
\hline
Hyperlink & ওয়েব/ইমেইল লিংক & Website \\
\hline
\end{tabular}
\end{center}

\begin{keypointbox}{ডেটা টাইপ নির্বাচন}
যে ফিল্ডে যোগ, বিয়োগ বা গড় বের করতে হবে সেটি Number/Currency হওয়া উচিত। যে নম্বরের শুরুতে 0 থাকতে পারে বা গণনা দরকার নেই, সেটি Text রাখা ভালো।
\end{keypointbox}

\begin{boardbox}{বোর্ড ফোকাস}
প্রশ্নে যদি ``মোবাইল নম্বরের জন্য কোন ডেটা টাইপ?'' আসে, উত্তর হবে Text/Short Text; কারণ মোবাইল নম্বর গাণিতিক হিসাবের জন্য নয় এবং শুরুতে 0 থাকতে পারে।
\end{boardbox}

% ৭. Query
\section{Query}

\subsection{Query কী}
\begin{itemize}
\item Database থেকে শর্তমত record বের করার পদ্ধতি
\end{itemize}

\subsection{Query-এর ধরন ও Condition}
\begin{itemize}
\item Select, Parameter, Action, Update, Delete, Append, Make Table
\item Condition: =, >, <, >=, <=, <>, AND, OR, LIKE, BETWEEN
\end{itemize}

\subsection{প্র্যাকটিক্যাল উদাহরণ}
\begin{itemize}
\item GPA 5.00 ছাত্রদের তালিকা, salary>20000 তালিকা ইত্যাদি
\end{itemize}

% ৮. Sorting
\section{Sorting}

\subsection{Sorting কী}
\begin{itemize}
\item Ascending ও Descending order-এ record সাজানো
\end{itemize}

\subsection{ব্যবহার}
\begin{itemize}
\item Alphabetical, roll/marks/salary/date অনুযায়ী সাজানো
\end{itemize}

% ৯. Indexing
\section{Indexing}
ইনডেক্সিং হলো বইয়ের সূচিপত্রের মতো একটি সহায়ক ব্যবস্থা, যা ডেটাবেজে নির্দিষ্ট রেকর্ড দ্রুত খুঁজে পেতে সাহায্য করে। বড় টেবিলে বারবার অনুসন্ধান করা হলে ইনডেক্স খুব কার্যকর।

\begin{definitionbox}{Indexing}
কোনো টেবিলের নির্দিষ্ট ফিল্ডের মানের ভিত্তিতে দ্রুত অনুসন্ধানের জন্য তৈরি বিশেষ ডেটা কাঠামোকে Index বলে; এই প্রক্রিয়াকে Indexing বলে।
\end{definitionbox}

\begin{keypointbox}{সুবিধা ও সীমাবদ্ধতা}
\begin{itemize}
    \item সুবিধা: অনুসন্ধান দ্রুত হয়, Sort ও Join অপারেশনে সহায়তা করে।
    \item সীমাবদ্ধতা: অতিরিক্ত স্টোরেজ লাগে এবং Insert/Update/Delete কিছুটা ধীর হতে পারে।
\end{itemize}
\end{keypointbox}

\begin{examplebox}{ব্যবহার}
StudentID, Roll, NID, AccountNo-এর মতো অনন্য বা বেশি অনুসন্ধানযোগ্য ফিল্ডে ইনডেক্স রাখা উপযোগী। কিন্তু Address বা Comment-এর মতো বড় লেখা ফিল্ডে সাধারণত ইনডেক্স কম দরকার।
\end{examplebox}

\begin{mistakebox}{Sorting বনাম Indexing}
Sorting ডেটা প্রদর্শনের ক্রম বদলায়। Indexing মূলত অনুসন্ধান দ্রুত করার জন্য সহায়ক কাঠামো তৈরি করে। দুইটি একই বিষয় নয়।
\end{mistakebox}

% ১০. Database Relation
\section{Database Relation}
রিলেশনাল ডেটাবেজে একাধিক টেবিল আলাদা রাখা হলেও প্রয়োজন অনুযায়ী তাদের মধ্যে সম্পর্ক তৈরি করা হয়। এই সম্পর্ক তৈরি হয় সাধারণত Primary key ও Foreign key-এর মাধ্যমে।

\subsection{Relation কী ও দরকারি ধরন}
\begin{definitionbox}{Primary Key ও Foreign Key}
\textbf{Primary Key} হলো এমন ফিল্ড যা প্রতিটি রেকর্ডকে অনন্যভাবে শনাক্ত করে। \textbf{Foreign Key} হলো এক টেবিলের এমন ফিল্ড, যা অন্য টেবিলের Primary key-এর সাথে সম্পর্ক তৈরি করে।
\end{definitionbox}

\begin{center}
\begin{tabular}{|p{0.24\textwidth}|p{0.45\textwidth}|p{0.20\textwidth}|}
\hline
রিলেশনের ধরন & ব্যাখ্যা & উদাহরণ \\
\hline
One-to-One & এক টেবিলের এক রেকর্ড অন্য টেবিলের এক রেকর্ডের সাথে যুক্ত & Person--Passport \\
\hline
One-to-Many & এক টেবিলের এক রেকর্ড অন্য টেবিলের একাধিক রেকর্ডের সাথে যুক্ত & Customer--Order \\
\hline
Many-to-Many & দুই পাশে একাধিক রেকর্ড যুক্ত; সাধারণত Junction table লাগে & Student--Course \\
\hline
\end{tabular}
\end{center}

\subsection{Junction table আলোচনা}
\begin{examplebox}{চিত্রের স্থানধারক}
Customer(CustomerID, Name) এবং Order(OrderID, CustomerID, Date) টেবিল দেখিয়ে CustomerID-এর মাধ্যমে One-to-Many relation আঁকা হবে।
\end{examplebox}

\begin{boardbox}{বোর্ড ফোকাস}
HSC পরীক্ষায় One-to-Many relation সবচেয়ে বেশি আসে। Parent table-এর Primary key child table-এ Foreign key হিসেবে ব্যবহার হয়।
\end{boardbox}

% ১১. Form
\section{Form}
Form হলো ডেটাবেজে ডেটা ইনপুট, দেখা ও সম্পাদনার জন্য ব্যবহারবান্ধব ইন্টারফেস। সরাসরি টেবিলে কাজ করলে ভুল হওয়ার সম্ভাবনা বেশি, তাই ব্যবহারকারীর জন্য ফর্ম তৈরি করা হয়।

\begin{definitionbox}{Form}
ডেটাবেজের টেবিল বা কুয়েরির ডেটা সহজভাবে এন্ট্রি, প্রদর্শন ও সম্পাদনার জন্য তৈরি ইন্টারফেসকে Form বলে।
\end{definitionbox}

\begin{keypointbox}{ফর্মের উপাদান}
Label, Text box, Combo box, List box, Check box, Radio button, Command button, Date picker ইত্যাদি কন্ট্রোল ফর্মে ব্যবহার করা যায়।
\end{keypointbox}

\begin{examplebox}{চিত্রের স্থানধারক}
Student Entry Form: StudentID, Name, GroupName, BirthDate, Mobile এবং Save/Search/Delete বোতামসহ একটি ফর্মের নকশা।
\end{examplebox}

\begin{boardbox}{ব্যবহারিক গুরুত্ব}
MS Access-এ টেবিলের উপর ভিত্তি করে Form Wizard বা Design View ব্যবহার করে ফর্ম তৈরি করা যায়। ফর্মের মাধ্যমে ডেটা এন্ট্রি করলে ব্যবহারকারীকে পুরো টেবিলের কাঠামো জানতে হয় না।
\end{boardbox}

% ১২. Report
\section{Report}

\subsection{Report কী}
\begin{itemize}
\item Query বা table থেকে print-friendly output তৈরি
\end{itemize}

\subsection{প্রকারভেদ ও ব্যবহার}
\begin{itemize}
\item Group report, summary report, student result, salary, sales reports
\end{itemize}

% ১৩. Corporate Database
\section{Corporate Database}

\subsection{সম্প্রসারণ ও ব্যবহার}
\begin{itemize}
\item Employee, Customer, Sales, Inventory, Accounting, HR, Payroll ইত্যাদি
\end{itemize}

\subsection{সফটওয়্যার ও সুবিধা}
\begin{itemize}
\item Oracle, DB2, SQL Server, Teradata ইত্যাদি; centralized data ও decision support
\end{itemize}

% ১৪. সরকারি প্রতিষ্ঠানে ডেটাবেজ
\section{সরকারি প্রতিষ্ঠানে ডেটাবেজ}
সরকারি ডেটাবেজ নাগরিক সেবা, প্রশাসনিক কাজ, নীতি নির্ধারণ এবং স্বচ্ছতা বৃদ্ধিতে গুরুত্বপূর্ণ ভূমিকা রাখে। ডিজিটাল বাংলাদেশ/স্মার্ট বাংলাদেশ ধারণায় সরকারি ডেটাবেজের ব্যবহার আরও বেড়েছে।

\subsection{ধারণা ও ব্যবহার}
\begin{keypointbox}{উদাহরণ}
\begin{itemize}
    \item জাতীয় পরিচয়পত্র ও ভোটার ডেটাবেজ।
    \item জন্ম-মৃত্যু নিবন্ধন ডেটাবেজ।
    \item পাসপোর্ট, ভিসা ও ইমিগ্রেশন ডেটাবেজ।
    \item ভূমি রেকর্ড ও নামজারি ডেটাবেজ।
    \item আয়কর, ভ্যাট ও রাজস্ব ডেটাবেজ।
    \item শিক্ষা বোর্ড, ফলাফল ও ভর্তি ডেটাবেজ।
    \item স্বাস্থ্য, হাসপাতাল ও টিকা ডেটাবেজ।
\end{itemize}
\end{keypointbox}

\subsection{সুবিধা}
\begin{examplebox}{ব্যবহার}
জাতীয় পরিচয়পত্র নম্বর ব্যবহার করে নাগরিকের পরিচয় যাচাই, ব্যাংক হিসাব খোলা, সিম নিবন্ধন, সরকারি ভাতা প্রদান ও অনলাইন সেবা দেওয়া সহজ হয়।
\end{examplebox}

\begin{mistakebox}{গোপনীয়তা}
সরকারি ডেটাবেজে ব্যক্তিগত তথ্য থাকে। তাই অনুমতি ছাড়া তথ্য প্রকাশ, দুর্বল পাসওয়ার্ড, অনিরাপদ শেয়ারিং এবং ব্যাকআপহীনতা বড় ঝুঁকি।
\end{mistakebox}

% ১৫. Data Security
\section{Data Security}

\subsection{Security কী ও প্রয়োজনীয়তা}
\begin{itemize}
\item Unauthorized access, hacking, malware, accidental deletion, hardware failure
\item Privacy ও business continuity রক্ষা করা
\end{itemize}

\subsection{Security পদ্ধতি}
\begin{itemize}
\item Strong password, access control, backup, antivirus, firewall, encryption, audit log
\end{itemize}

% ১৬. Data Encryption
\section{Data Encryption}

\subsection{Encryption ও Decryption}
\begin{itemize}
\item Encryption কী, Decryption কী, প্রয়োজনীয়তা
\end{itemize}

\subsection{ধরন ও ব্যবহার}
\begin{itemize}
\item Symmetric ও Asymmetric, Public/Private key, online banking, e-commerce
\end{itemize}

% ১৭. Practical Database Task
\section{Practical Database Task}
এই অধ্যায়ের ব্যবহারিক অংশে সাধারণত MS Access বা সমজাতীয় DBMS ব্যবহার করে ডেটাবেজ, টেবিল, কুয়েরি, ফর্ম ও রিপোর্ট তৈরি করতে বলা হয়।

\begin{practicebox}
\textbf{কাজ: StudentDB তৈরি}
\begin{enumerate}
    \item StudentDB নামে একটি ডেটাবেজ তৈরি করো।
    \item Student নামে একটি টেবিল তৈরি করো।
    \item ফিল্ড দাও: StudentID, Roll, Name, GroupName, GPA, Mobile।
    \item StudentID-কে Primary Key করো।
    \item অন্তত ৫টি রেকর্ড এন্ট্রি করো।
    \item GPA 5.00 পাওয়া ছাত্রদের দেখানোর জন্য Select query তৈরি করো।
    \item Name অনুযায়ী Ascending sort করো।
    \item Student Entry Form এবং GPA Report তৈরি করো।
\end{enumerate}
\end{practicebox}

\begin{examplebox}{নমুনা কুয়েরি}
\begin{verbatim}
SELECT Roll, Name, GPA
FROM Student
WHERE GPA = 5.00
ORDER BY Name ASC;
\end{verbatim}
\end{examplebox}

\begin{keypointbox}{ব্যবহারিক খাতার জন্য}
প্রতিটি কাজের শেষে স্ক্রিনশট/চিত্রের স্থানধারক রাখা যেতে পারে: Table Design View, Datasheet View, Query Design, Form View, Report Preview।
\end{keypointbox}

% ১৮. অনুশীলনী
\section{অনুশীলনী}

\begin{itemize}
\item MCQ, জ্ঞানমূলক, অনুধাবনমূলক, প্রয়োগমূলক, Database design ও Query লিখন
\end{itemize}


% Resource linked: resources/HSC ICT Chapter 06 Handnote.pdf
